%% annex/p07/calculi-and-cut-elimination.tex --- evidence for P7.1--7.2

\section{Calculus identities and Cut elimination}

\begin{annexitem}{annex:p7.calculi.1}{thm:p07-lkd-soundness}{Soundness of the defined-negation sequent calculus}
\sketchnote{the individual semantic calculations for every structural,
propositional, and quantifier rule; in particular, the substitution-lemma
calculation and the two eigenvariable cases.}
Induct on the height of a
\(\mathsf{LK}_{\mathrm d}\)-derivation.  The initial and structural rules
preserve the semantic reading of the sequent.  Each propositional case unfolds
the truth clause of its principal connective.  The quantifier cases use the
first-order substitution lemma; freshness of the eigenvariable permits the
assignment variation in \(\forall R\) and \(\exists L\).  Cut composes the
two validity implications through its displayed formula.
\end{annexitem}

\begin{annexitem}{annex:p7.calculi.2}{thm:p07-lk-presentation-equivalence}{Canonical/defined-negation LK equivalence}
\sketchnote{contexting uncontexted identity and its converse; admissibility
or derivability of the translated primitive-negation rules; both directions
between shared-context and split-context implication-left using weakening and
contraction; the remaining propositional cases; and all four
quantifier/eigenvariable cases.  No cut-free, height, or size preservation is
established.}
Translate formulas occurrencewise by \(d\).  In the canonical-to-defined
direction, induct on the canonical derivation: simulate contexted initial
sequents by uncontexted identity plus weakenings, expand primitive negation as
implication to \(\bot\), and split shared implication-left contexts using the
registered structural rules.  In the reverse direction, use canonical
contexted identity and structural rules to simulate the uncontexted and
split-context presentation.  Homomorphic connective and quantifier clauses
then preserve each translated endsequent.  The omitted rule-by-rule
verification is why the result is a sketch.
\end{annexitem}

\begin{annexitem}{annex:p7.cut.1}{thm:p07-cut-mix-equivalence}{Equipollence of Cut and Mix}
\begin{proof}
Cut is Mix with both multiplicities one.  Conversely, for multiplicities
\(m,n\), use \(m-1\) CR steps on the left premise and \(n-1\) CL steps on the
right, then apply Cut.  In LJ the succedent already has one copy, so only CL
is needed.  Replacing every Mix instance preserves its endsequent.  Thus the
formulations derive the same sequents, and eliminating Mix after reading Cuts
as unit Mixes eliminates Cut.
\end{proof}
\end{annexitem}

\begin{annexitem}{annex:p7.cut.2}{thm:p07-cut-elimination}{Cut elimination for defined-negation LK and LJ}
\sketchnote{the complete calculus-specific table of commuting and
principal reductions for every rule of \(\mathsf{LK}_{\mathrm d}\) and
\(\mathsf{LJ}\), including every contraction-multiplicity calculation and
all quantifier/eigenvariable cases.}
By Annex~\ref{annex:p7.cut.1}, read every Cut as a unit-multiplicity Mix.
Choose an uppermost Mix of maximal
\(\operatorname{rk}_{\mathrm{fm}}\).  If its formula is nonprincipal,
commute it above the relevant final inference, lowering
\(\operatorname{ht}\).
Delete initial and weakening cases.  When contraction is principal, expose
the uncontracted premise and increase the relevant multiplicity, again
lowering height.  When both occurrences are principal, replace the Mix by
Mixes on proper subformulas, together with any residual same-rank Mixes of
smaller height.  Before a quantifier permutation, alpha-rename the
eigenvariable and apply capture-avoiding substitution.  The lexicographic
\(\operatorname{rk}_{\mathrm{fm}}\)--\(\operatorname{ht}\) measure is well
founded, so the omitted reduction table would
iterate to a Mix-free, hence cut-free, derivation of the same endsequent.
\end{annexitem}

\begin{annexitem}{annex:p7.cut.3}{thm:p07-cut-free-subformula}{Cut-free propositional subformula property}
\begin{proof}
Induct on derivation height: every premise formula is copied from the
conclusion or is an immediate subformula of its principal formula.
\end{proof}
\end{annexitem}

\begin{annexitem}{annex:p7.cut.4}{thm:p07-propositional-analytic-existence}{Analytic existence for derivable propositional sequents}
\sketchnote{the complete Cut-elimination reduction table omitted from
Annex~\ref{annex:p7.cut.2}; the subformula step itself is proved in
Annex~\ref{annex:p7.cut.3}.}
Apply the sketched Cut-elimination transformation to the supplied derivation.
The result has the same endsequent and no Cut.  Annex~\ref{annex:p7.cut.3}
then gives the literal subformula property, which is exactly analyticity for
the propositional active calculi.
\end{annexitem}

\begin{annexitem}{annex:p7.cut.5}{thm:p07-empty-sequent}{Underivability of the empty sequent}
\sketchnote{the complete Cut-elimination reduction table on which the passage
to a cut-free derivation depends.}
If the empty sequent were derivable, the cut-elimination sketch would give a
cut-free derivation of it.  It is not an instance of either zero-premise rule,
and the conclusion of every other cut-free rule on
\texttt{sheet:lk-defneg} contains a displayed formula occurrence.  Hence no
cut-free derivation has the empty sequent as its root.
\end{annexitem}
